{

An object of this class represents the original executable or a library. It serves as a container of BPatch_module

objects.}



{

std::string name()}



{

std::string pathName()}



{

Return the name of this file; either just the file name or the fully path-qualified name.



{

Dyninst::Address fileOffsetToAddr(Dyninst::Offset offset)}



{

Convert the provided offset into the file into a full address in memory.



{

struct Region {}



{

typedef enum { UNKNOWN, CODE, DATA } type_t;}



{

Dyninst::Address base;}



{

unsigned long size;}



{

type_t type;}



{

};}



{

void regions(std::vector<Region> &regions)}



{

Returns information about the address ranges occupied by this object in memory.



{

void modules(std::vector<BPatch_module *> &modules)}



{

Returns the modules contained in this object.}



{

std::vector<BPatch_function*> *findFunction\index{findFunction: : : : : : :!}(}



{

const char *name,}



{

std::vector<BPatch_function*> &funcs,}



{

bool showError = true,}



{

bool regex_case_sensitive = true,}



{

bool incUninstrumentable = false)}



{

Return a vector of BPatch_functions corresponding to name, or NULL if the function does not exist. If name contains a

POSIX-extended regular expression, and dont_use_regex is false, a regular expression search will be performed on

function names and matching BPatch_functions returned. If showError is true, then Dyninst will report an error via

the BPatch::registerErrorCallback if no function is found.}



{

If the incUninstrumentable flag is set, the returned table of procedures will include uninstrumentable functions. The

default behavior is to omit these functions.}



{

[NOTE: If name is not found to match any demangled function names in the module, the search is repeated as if name is a

mangled function name. If this second search succeeds, functions with mangled names matching name are returned

instead.]}



{

bool findPoints(Dyninst::Address addr,



{

std::vector<BPatch_point *> &points);}



{

Return a vector of BPatch_points that correspond with the provided address, one per function that includes an

instruction at that address. There will be one element if there is not overlapping code.}



{

std::vector<BPatch_function*> *findFunction\index{findFunction: : : : : : :!}(}



{

const char *name,}



{

std::vector<BPatch_function*> &funcs,}



{

bool notify_on_failure = true,}



{

bool regex_case_sensitive = true,}



{

bool incUninstrumentable = false)}



{

Return a vector of BPatch_functions matching name, or NULL if the function does not exist. If name contains a

POSIX-extended regular expression, a regex search will be performed on function names, and matching BPatch_functions

returned. [NOTE: The std::vector argument funcs must be declared fully by the user before calling this function.

Passing in an uninitialized reference will result in undefined behavior.]}



{

If the incUninstrumentable flag is set, the returned table of procedures will include uninstrumentable functions. The

default behavior is to omit these functions.}



{

[NOTE: If name is not found to match any demangled function names in the BPatch_object, the search is repeated as if

name is a mangled function name. If this second search succeeds, functions with mangled names matching name are

returned instead.]}





\bigskip



